\documentclass[11pt]{article}
\usepackage{hyperref}
\title{A designation can expire too early even after the return settles}
\date{}
\begin{document}\maketitle

\paragraph{Scope.}
With the \href{https://github.com/eternal-november/eos-research-programmes/pull/582}
{frozen deferred-placement receiver}, a true local ingress observation
selects the alternate bay for an old compulsory shelf return. The earlier
proof requires the demand to stay consistent with the issued choice.
The present test asks if this requirement can end when the old return
actually settles and the carrier has cleared. It cannot in the modeled
family. This is not a claim about every reservation mechanism.

\paragraph{Discriminator.}
Keep the old five-tick unload deadline and give the next carrier a
release at tick~8 followed by a due goal at tick~9. The selected stage
step at tick~1 physically identifies the first demand. The controller
unloads a standing shelf in the other bay and clears it by tick~6.
The source schedule at tick~7 may then legitimately reassign the new due
task to the other staged route, without altering the old command or
using the previous, still-staged carrier. The final selected robot takes
a real native ingress step and then attempts its two native forward moves.
There is enough time for this revised job as such.

Three matched checks distinguish mechanisms. If the first assignment
remains, the unchanged receiver's old return and subsequent loaded
delivery both succeed. If the schedule reassigns after the old return
and clearance, the placement is now in the required bay and the new
delivery fails. If the return instead physically deposited its shelf
in the other permissible bay, the identical final demand, ingress step,
deadline, resources and unrelated service succeed. This last comparison
fixes a placement with hindsight to isolate its physical effect; it is
not a deployable policy that knows the later assignment at the first cut.
Both service statements count native delivery rather than mere
nominal access, and independent service is unchanged throughout.
The early report in the failed branch is true, not delayed or forged.

\paragraph{Exact dependency at this scope.}
For one standing old shelf that cannot be removed before the single
loaded crossing, the next route succeeds exactly when its required
cell differs from that shelf's cell, with the earlier standing shelf and
the timing conditions already satisfied. If a new authority may select
any cell in a set of possible final requirements, robust service requires
the actual occupied cell to be outside this set. The positive proof
must protect that disjointness until the last affected passage, not just
until signed acceptance, physical unload, or the carrier's departure.
A timely removal or a genuinely admissible alternate route would be
a different way to discharge the dependency. Merely rejecting the
new request may protect safety but does not serve the new obligation.

The inherited source and policy are pinned. For all six blocked-bay
and reassigned-demand combinations the old and disjoint deliveries
are preserved, and the stable, changed, and alternate-footprint runs
agree with the direct standing-cell condition. These 18 authored native
checks illustrate an effect lifetime, not an independent mission,
operational dominance, or an argument for extending any real-world
authorization by default.
\end{document}
